% Options for packages loaded elsewhere
\PassOptionsToPackage{unicode}{hyperref}
\PassOptionsToPackage{hyphens}{url}
\PassOptionsToPackage{dvipsnames,svgnames,x11names}{xcolor}
\PassOptionsToPackage{space}{xeCJK}
\documentclass[
  10pt,
]{article}
\usepackage{xcolor}
\usepackage[margin=2.2cm]{geometry}
\usepackage{amsmath,amssymb}
\setcounter{secnumdepth}{-\maxdimen} % remove section numbering
\usepackage{iftex}
\ifPDFTeX
  \usepackage[T1]{fontenc}
  \usepackage[utf8]{inputenc}
  \usepackage{textcomp} % provide euro and other symbols
\else % if luatex or xetex
  \usepackage{unicode-math} % this also loads fontspec
  \defaultfontfeatures{Scale=MatchLowercase}
  \defaultfontfeatures[\rmfamily]{Ligatures=TeX,Scale=1}
\fi
\usepackage{lmodern}
\ifPDFTeX\else
  % xetex/luatex font selection
  \setmainfont[]{Times New Roman}
  \ifXeTeX
    \usepackage{xeCJK}
    \setCJKmainfont[]{Songti SC}
  \fi
  \ifLuaTeX
    \usepackage[]{luatexja-fontspec}
    \setmainjfont[]{Songti SC}
  \fi
\fi
% Use upquote if available, for straight quotes in verbatim environments
\IfFileExists{upquote.sty}{\usepackage{upquote}}{}
\IfFileExists{microtype.sty}{% use microtype if available
  \usepackage[]{microtype}
  \UseMicrotypeSet[protrusion]{basicmath} % disable protrusion for tt fonts
}{}
\makeatletter
\@ifundefined{KOMAClassName}{% if non-KOMA class
  \IfFileExists{parskip.sty}{%
    \usepackage{parskip}
  }{% else
    \setlength{\parindent}{0pt}
    \setlength{\parskip}{6pt plus 2pt minus 1pt}}
}{% if KOMA class
  \KOMAoptions{parskip=half}}
\makeatother
\usepackage{longtable,booktabs,array}
\usepackage{caption}
\captionsetup[table]{skip=6pt}
\newcounter{none} % for unnumbered tables
\usepackage{calc} % for calculating minipage widths
% Correct order of tables after \paragraph or \subparagraph
\usepackage{etoolbox}
\makeatletter
\patchcmd\longtable{\par}{\if@noskipsec\mbox{}\fi\par}{}{}
\makeatother
% Allow footnotes in longtable head/foot
\IfFileExists{footnotehyper.sty}{\usepackage{footnotehyper}}{\usepackage{footnote}}
\makesavenoteenv{longtable}
\setlength{\emergencystretch}{3em} % prevent overfull lines
\providecommand{\tightlist}{%
  \setlength{\itemsep}{0pt}\setlength{\parskip}{0pt}}
\usepackage{bookmark}
\IfFileExists{xurl.sty}{\usepackage{xurl}}{} % add URL line breaks if available
\urlstyle{same}
% fallback for those not using the hyperref driver hyperxmp:
\makeatletter
\@ifundefined{xmpquote}{\newcommand{\xmpquote}[1]{#1}}{}
\makeatother
\hypersetup{
  pdftitle={Interface First{,} Then Teach{,} Then Consolidate: Self-Evolution of a 4B Local Computer-Use Agent on a Four-Dollar Teacher Budget},
  pdfauthor={Wenbo Yan},
  colorlinks=true,
  linkcolor={blue},
  filecolor={Maroon},
  citecolor={Blue},
  urlcolor={blue},
  pdfcreator={LaTeX via pandoc}}

\title{Interface First, Then Teach, Then Consolidate: Self-Evolution of
a 4B Local Computer-Use Agent on a Four-Dollar Teacher Budget}
\author{Wenbo Yan}
\date{October 2026}

\begin{document}
\maketitle


\emph{Draft v0.3, 2026-10-26 (§7 daily use rewritten: WeChat,
decision-model routing, latency; references added). Every number is
taken from a file under \texttt{evidence/} (cited in brackets).
Concurrent work (Terminal-RSI, Harness²) is discussed in §7. Statistical
tests are two-sided Fisher exact tests over task runs unless stated;
small samples are marked as directional.}

\subsection{Abstract}\label{abstract}

We study how far a 4-billion-parameter local model (Agents-A1-4B, 4-bit,
llama.cpp) can improve itself at computer-use tasks --- file editing,
web forms and native macOS applications --- when a cloud teacher is
available but capped at US\$4 for the whole project (US\$1.83 spent).
Four findings come from controlled same-session comparisons on frozen
benchmarks with programmatic checkers; a fifth from putting the agent in
daily use:

\begin{enumerate}
\def\labelenumi{\arabic{enumi}.}
\tightlist
\item
  \textbf{Learning from demonstrations works in context, on file tasks.}
  Trajectory memory or verified skills raise in-family success from
  14/48 to 39--42/48 after a 60-task stream; skills are not better than
  plain memory.
\item
  \textbf{On web and desktop tasks, small students retrieve guidance but
  do not follow it.} Correct teacher guidance is present in the prompt,
  yet the student's actions do not change. Before calling this a
  capability gap, we found that several ``capability'' failures were
  \emph{interface} failures: making one observation field explicit takes
  a settings task from 1/18 to 17/18 (p \textless{} 1e-5); replacing
  free-text editing with line-based tools takes desktop TextEdit tasks
  from 4/18 to 17/18 (p = 2e-5).
\item
  \textbf{Teacher lessons are adopted only when the student must ground
  them in the current task.} On unseen desktop tasks, lessons in the
  system prompt change nothing (3/15 vs 3/15); asking the student to
  first write how each lesson applies to this task gives 10/15 (p =
  0.025), and a plan-only control separates the two effects.
\item
  \textbf{Consolidating guided successes into LoRA weights removes the
  need for guidance in-family (30/72 → 65/72) but erodes held-out
  code-writing (15/36 → 1/36).} Across five ablations, gain and
  forgetting come from the \emph{same} training samples; replay (random
  or gradient-selected) protects only capabilities its data covers, and
  the trained student stops verifying its own output (48/72 → 0).
\item
  \textbf{In daily use, the 4B model should not be the one making
  choices.} A hosted decision model answering one multiple-choice
  question per request routes it to a program handler: 30/30 on two
  held-out phrasing sets where hand-written patterns get 15/30 and
  16/30. Measured per stage, the model's latency is prompt reading and
  long retelling, not reasoning: routing a weather question away from
  web search takes it from 36 s to 0.8 s, and a knowledge question from
  16 s (and a wrong answer) to 2.3 s.
\end{enumerate}

With the resulting interface and grounded lessons, the student solves
41/45 native macOS tasks through the Accessibility API (baseline 19/30
on the training tasks, 3/15 on unseen variants). We argue that a
self-improving small agent should, in order: rule out interface causes,
teach procedures in a form the student adopts, and only then consolidate
--- while monitoring held-out behaviour, not just in-family scores.

\subsection{1. Introduction}\label{introduction}

Small local models are private, cheap and always available, but weak. A
useful personal agent built on one should improve with use: notice
failures, get help at bounded cost, keep what it learns, and not apply
it where it does not belong. Recent self-improvement work either changes
the \emph{harness} around a strong model (prompts, skills, memory;
Harness², Meta-Harness) or changes the \emph{weights} of a mid-size
model with large compute (Terminal-RSI, 27B). Both implicitly assume the
model follows the guidance it is given.

We test that assumption at the other end of the scale: a 4B model on
consumer hardware (Kaggle T4, an Apple M4 laptop), a teacher (DeepSeek)
behind a global US\$4 ledger, and three task domains with hidden
checkers. Our contribution is less a new method than a set of controlled
findings about \emph{where} a small self-improving agent fails and what
fixes each failure:

\begin{itemize}
\tightlist
\item
  an \textbf{adoption gap}: retrieval is reliable, adoption is not (§4);
\item
  \textbf{interface artefacts} that masquerade as capability gaps, found
  with one-field ablations (§5);
\item
  \textbf{grounded lessons}: a cheap prompting change that makes teacher
  rules effective on unseen tasks (§5.3);
\item
  \textbf{consolidation trade-offs}: a five-version ablation that
  localises forgetting to the samples that carry the gain (§6);
\item
  \textbf{choices to a decision model, actions to the program} in daily
  use, with the failure modes that motivated it (invented tool results,
  look-alike requests) and a per-stage latency analysis (§7).
\end{itemize}

All code, frozen benchmarks, raw runs and per-experiment reports are
public; negative results are reported.

\subsection{2. System}\label{system}

\textbf{Student and loop.} Agents-A1-4B (Q4\_K\_M) served by llama.cpp;
an agentic loop with at most 5--8 steps and 6--10 model calls per task;
tools depend on the domain (workspace file tools; browser observe/act on
DOM snapshots; macOS Accessibility press/fill/menu). The checker never
reaches the agent.

\textbf{Learning loop.} After a failure the teacher solves the task from
a clean state (cached per task instance); a verified demonstration
becomes trajectory memory and/or a parameterised, sandbox-verified
skill. A TF-IDF relevance gate shows learned items only for tasks close
to where they were verified; when nothing passes, the request is
byte-identical to the raw student's.

\textbf{Lessons.} After a failure the teacher compares the student's
actions with the checker's finding and writes one or two general
imperative rules without task-specific values; rules are gated like
memory.

\textbf{Consolidation loop (R).} Following RSR-style re-execution: the
student re-attempts practice variants \emph{with} guidance
(demonstration + rules), only verified successes are kept, the guidance
is removed, and the student is fine-tuned with LoRA (rank 16 on
attention-output and MLP projections, response-only loss) and evaluated
without guidance.

\textbf{Decision model.} For choices rather than text, a hosted decision
model (Cloudflare Clef, 27B) answers typed questions --- yes/no or one
of named options --- with calibrated probabilities in about 0.3 s. It
never writes text or acts. It extracts per-item goals for a web checker
(§5.1) and routes daily requests (§7.2).

\textbf{Budget.} Every cloud call reserves against a global ledger
before it is sent (US\$4 total; US\$1.83 spent over all experiments,
including void runs; decision-model calls are booked at list price,
US\$0.10 in total, although they ran inside the provider's free
allocation) {[}ledger, \texttt{core/evolution/teacher\_budget.py}{]}.

\subsection{3. Benchmarks}\label{benchmarks}

\begin{itemize}
\tightlist
\item
  \textbf{Kage-Stream (files)}: 4 learned families, 2 held-out families
  (text normalisation, config migration) and 3 read-only negative
  families; procedurally generated, frozen by hash, calibrated so the
  raw student passes 10--50\% {[}\texttt{eval/stream\_v1}{]}.
\item
  \textbf{Kage-Web}: local fixtures for invoices (hidden rows),
  notification settings, hotel catalogue (ratings behind ``Details''), a
  two-step profile form (held-out) and read-only lookups
  {[}\texttt{eval/stream\_web\_v1}{]}.
\item
  \textbf{Desktop-v1 (macOS)}: Calculator (4 + 2 unseen chain variants),
  TextEdit (3 + 3 unseen variants, dedicated files) and read-only System
  Settings (3), driven only through the Accessibility API --- no
  keyboard events, no coordinates; checkers read the app or the disk
  themselves {[}\texttt{eval/desktop\_v1}{]}.
\end{itemize}

\textbf{Evaluation discipline.} GPU inference is not bit-reproducible
even at temperature 0: four identical raw-student evaluations gave
7--13/24 in-family {[}\texttt{docs/report/kage-tech-report.md} §5.4{]}.
We therefore compare arms in the same session with ≥ 3 runs and treat
differences inside that range as noise.

\subsection{4. Learning in context, and the adoption
gap}\label{learning-in-context-and-the-adoption-gap}

\textbf{Files (E-1).} After a 60-task stream, trajectory memory reaches
42/48 and a skill library 39/48 in-family against 14/48 for the raw
student and 12/48 for a budget-matched retry baseline (skills vs raw: p
= 6e-8; skills vs memory: 5 vs 8 discordant pairs, p = 0.58) for
US\$0.07--0.20 per arm. Learned knowledge distracted the student on
held-out families until the relevance gate restored them to the raw
student's range with zero misuse {[}tech report §5{]}.

\textbf{Web (C2, C3).} The same memory barely helps on web tasks
(learned arms 20/72 vs raw 12/54, p = 0.54). Per-task inspection shows
retrieval is not the problem: every evaluation prompt contained the
relevant demonstration, yet the student followed it in 1--2 of 6 invoice
tasks {[}\texttt{evidence/reports/c2}{]}. Teacher-written rules fix a
\emph{missing-step} failure (invoices: 0/24 → 14/24 with rules)
{[}\texttt{evidence/reports/c3}{]}; guided re-execution also helps only
where a step or idiom is missing (guided 71/112 vs unguided 34/112, p =
1e-6, but settings 6 vs 6 and hotels 7 vs 8)
{[}\texttt{evidence/reports/r2}{]}.

\subsection{5. Interface first}\label{interface-first}

That the interface between agent and computer matters is not new:
SWE-agent designed an agent-computer interface (file viewer, editor with
linting, search commands) and showed large gains for
software-engineering agents built on strong models (Yang et al., 2024).
Our contribution here is narrower and diagnostic: with a 4B student we
isolate \emph{single} interface factors --- one observation field, one
tool family --- in controlled ablations, and show that failures we had
attributed to missing capability, and that demonstrations, rules and
consolidation did not fix, disappear when that one factor changes.

\subsubsection{5.1 Web: one field (C4)}\label{web-one-field-c4}

Settings tasks resisted demonstrations, rules and consolidation, and we
had cited them as ``knows but cannot do''. Our compact observation
omitted \texttt{checked:\ false} for unchecked boxes. With the raw
student and no guidance {[}\texttt{evidence/reports/c4/analysis.md}{]}:

{\def\LTcaptype{none} % do not increment counter
\begin{longtable}[]{@{}llll@{}}
\toprule\noalign{}
Interface & Settings (/18) & Total (/90) & Max input tokens \\
\midrule\noalign{}
\endhead
\bottomrule\noalign{}
\endlastfoot
v1 (unchecked boxes have no \texttt{checked} key) & 1 & 29 & 7,545 \\
v2 (explicit \texttt{checked:\ true/false}) & 14 & 41 & 7,406 \\
v2 + context packing + state table & 17 & 39 & 3,946 \\
\end{longtable}
}

A second round {[}\texttt{evidence/reports/c4/round2.md}{]} re-baselined
the benchmark on v2 (41/90, settings 18/18). A pre-save checker --- Clef
extracts per-item goals (48/48 correct), a program compares them with
the page --- never fired (0 of 59 checks): once the interface was fixed
there was nothing left for it to catch. Hotel failures, which we first
attributed to context packing, are procedural: the student never opens
``Details'' to compare ratings, with or without packing (7 vs 4/18, p =
0.47).

\subsubsection{5.2 Desktop: line-based editing
(G0)}\label{desktop-line-based-editing-g0}

Driving TextEdit through the Accessibility API with a
\texttt{fill(text)} tool, the student fails 14 of 18 edits, all on line
breaks: it drops the final newline or writes a literal
\texttt{\textbackslash{}n}. Teacher rules plus planning do not help
(4/12) --- the student \emph{states} that each line ends with a newline
and still emits the wrong JSON string. Tools that remove the need to
encode control characters solve it
{[}\texttt{evidence/reports/g0/analysis.md} §5{]}:

{\def\LTcaptype{none} % do not increment counter
\begin{longtable}[]{@{}
  >{\raggedright\arraybackslash}p{(\linewidth - 8\tabcolsep) * \real{0.2000}}
  >{\raggedright\arraybackslash}p{(\linewidth - 8\tabcolsep) * \real{0.2000}}
  >{\raggedright\arraybackslash}p{(\linewidth - 8\tabcolsep) * \real{0.2000}}
  >{\raggedright\arraybackslash}p{(\linewidth - 8\tabcolsep) * \real{0.2000}}
  >{\raggedright\arraybackslash}p{(\linewidth - 8\tabcolsep) * \real{0.2000}}@{}}
\toprule\noalign{}
\begin{minipage}[b]{\linewidth}\raggedright
TextEdit tools
\end{minipage} & \begin{minipage}[b]{\linewidth}\raggedright
Fix typo (/6)
\end{minipage} & \begin{minipage}[b]{\linewidth}\raggedright
Append line (/6)
\end{minipage} & \begin{minipage}[b]{\linewidth}\raggedright
Write lines (/6)
\end{minipage} & \begin{minipage}[b]{\linewidth}\raggedright
Total
\end{minipage} \\
\midrule\noalign{}
\endhead
\bottomrule\noalign{}
\endlastfoot
\texttt{fill} (whole text) & 0 & 3 & 1 & 4/18 \\
+ \texttt{replace\_text}, \texttt{append\_text} & 6 & 2 & 0 & 8/18 \\
\texttt{replace\_text} + \texttt{append\_lines} / \texttt{set\_lines} &
6 & 5 & 6 & \textbf{17/18} \\
\end{longtable}
}

Two smaller interface slips on the desktop (the student passed a
control's id or the whole observation line as its reference) were fixed
by accepting both.

\subsubsection{5.3 Grounded lessons (G0)}\label{grounded-lessons-g0}

Calculator chain tasks fail on operator precedence: the student presses
\texttt{15\ +\ 27\ ×\ 3\ =} and gets 96. The teacher wrote one rule from
the failure (``when an expression mixes addition and multiplication,
press Equals after the addition\ldots{}'', cost US\$0.0005 for all four
rules). On five unseen variants (three runs each):

{\def\LTcaptype{none} % do not increment counter
\begin{longtable}[]{@{}
  >{\raggedright\arraybackslash}p{(\linewidth - 6\tabcolsep) * \real{0.2500}}
  >{\raggedright\arraybackslash}p{(\linewidth - 6\tabcolsep) * \real{0.2500}}
  >{\raggedright\arraybackslash}p{(\linewidth - 6\tabcolsep) * \real{0.2500}}
  >{\raggedright\arraybackslash}p{(\linewidth - 6\tabcolsep) * \real{0.2500}}@{}}
\toprule\noalign{}
\begin{minipage}[b]{\linewidth}\raggedright
Arm
\end{minipage} & \begin{minipage}[b]{\linewidth}\raggedright
Calculator chains (/6)
\end{minipage} & \begin{minipage}[b]{\linewidth}\raggedright
TextEdit (/9)
\end{minipage} & \begin{minipage}[b]{\linewidth}\raggedright
Total (/15)
\end{minipage} \\
\midrule\noalign{}
\endhead
\bottomrule\noalign{}
\endlastfoot
no lessons & 1 & 2 & 3 \\
lessons in the system prompt & 1 & 2 & \textbf{3} \\
``write a plan first'', no lessons & 4 & 2 & 6 \\
lessons next to the task + ``say how each lesson applies, then act'' & 5
& 5 & \textbf{10} (p = 0.025) \\
\end{longtable}
}

In the system prompt the rules are ignored: the student still presses
\texttt{9\ +\ 6\ ×\ 4\ =}. Requiring it to ground each rule in the
current task (an exact button order) makes it adopt them, and it
generalises the addition rule to a subtraction/division variant (``press
Equals to get 36, then divide''). The plan-only control shows that
planning alone explains most of the Calculator gain, while the TextEdit
gain needs the rule content.

\textbf{Final desktop configuration} (line-based tools, interface fixes,
grounded lessons; 15 tasks × 3 runs): \textbf{41/45} --- training tasks
29/30 (baseline 19/30, p = 0.003), unseen variants 12/15 (3/15, p =
0.003); Calculator 17/18, TextEdit 15/18, System Settings (read-only)
9/9 {[}\texttt{evidence/reports/g0/analysis.md} §6{]}.

\subsection{6. Consolidation: gain and forgetting come from the same
samples}\label{consolidation-gain-and-forgetting-come-from-the-same-samples}

R2 collected 1,900 guided successes (guidance stripped); R3 trains a
LoRA adapter and evaluates without guidance (file sets: 24 in-family +
12 held-out per run, 3 runs; web: 18 learned per run)
{[}\texttt{evidence/reports/r3*}{]}:

{\def\LTcaptype{none} % do not increment counter
\begin{longtable}[]{@{}
  >{\raggedright\arraybackslash}p{(\linewidth - 8\tabcolsep) * \real{0.2000}}
  >{\raggedright\arraybackslash}p{(\linewidth - 8\tabcolsep) * \real{0.2000}}
  >{\raggedright\arraybackslash}p{(\linewidth - 8\tabcolsep) * \real{0.2000}}
  >{\raggedright\arraybackslash}p{(\linewidth - 8\tabcolsep) * \real{0.2000}}
  >{\raggedright\arraybackslash}p{(\linewidth - 8\tabcolsep) * \real{0.2000}}@{}}
\toprule\noalign{}
\begin{minipage}[b]{\linewidth}\raggedright
Version
\end{minipage} & \begin{minipage}[b]{\linewidth}\raggedright
Training change
\end{minipage} & \begin{minipage}[b]{\linewidth}\raggedright
Files in-family (/72)
\end{minipage} & \begin{minipage}[b]{\linewidth}\raggedright
Held-out (/36)
\end{minipage} & \begin{minipage}[b]{\linewidth}\raggedright
Web learned (/54)
\end{minipage} \\
\midrule\noalign{}
\endhead
\bottomrule\noalign{}
\endlastfoot
raw student & --- & 30 & 15 (pooled 27/60) & 10 \\
R3 & guided successes & 64 (p = 2.6e-9) & \textbf{1} & 21 \\
R3b & drop reasoning-free targets + random replay & 34 & \textbf{15} &
31 \\
R3c & drop + gradient-selected (LESS) replay & 34 & 2 & 22 \\
R3d & teacher rationales for reasoning-free targets + LESS replay & 62 &
6 & 20 \\
R3e & rationales + random replay & \textbf{65} & \textbf{1} & 22 \\
\end{longtable}
}

General-ability check (30 questions): 27--28/30 for every version. Under
the round-2 web interface, R3e solves 7/18 invoices where the raw
student solves 0/54 (p = 2e-5): consolidated knowledge survives an
interface change.

LoRA is known to forget less than full fine-tuning, and also to learn
less (Biderman et al., 2024). Even so, a rank-16 adapter trained on
about 1,900 samples erased held-out code writing here; what the ablation
adds is \emph{which} samples do it. What the five versions show:

\begin{itemize}
\tightlist
\item
  \textbf{Gain and forgetting move together.} Held-out success collapses
  exactly when the \textasciitilde300 file-solution steps
  (\texttt{write\_file} targets, mostly without reasoning in the raw
  data) are trained (R3, R3d, R3e) and is preserved only when they are
  removed (R3b) --- which also removes the in-family gain. Replay type
  and the presence of reasoning do not change this.
\item
  \textbf{Replay protects only what it covers.} Our replay data
  (desktop-action text from AgentNet, MCP tool calls from Toucan)
  contains no code writing; neither random nor gradient-selected replay
  protected held-out \emph{code} tasks. Gradient-influence selection
  (LESS) picks data aligned with the target --- the opposite of what
  anti-forgetting replay needs.
\item
  \textbf{The mechanism is code quality plus a lost habit.} On held-out
  tasks the trained student writes plausible but wrong code
  (e.g.~\texttt{encode("ascii")} that raises on non-ASCII input) and no
  longer reads its output back: read-back verification after writing
  drops from 48/72 runs (raw) to 0 for every trained version, because
  training trajectories end as soon as the checker passes. ``Not writing
  reasoning'', our first explanation, was a symptom: R3d/R3e reason
  before every call and still forget.
\end{itemize}

Implications we did not test (compute ended with this study):
self-distilled replay of the raw student's own correct code on other
tasks, keeping write-then-verify steps in training trajectories, and
gating the adapter by task family at deployment.

\subsection{7. From benchmark to daily
use}\label{from-benchmark-to-daily-use}

\subsubsection{7.1 The program supplies information and
acts}\label{the-program-supplies-information-and-acts}

We then put the student behind the user's own requests on their Mac
(calendar, mail, open documents, apps, websites, WeChat) and logged
every turn. The benchmark pattern held in a sharper form: the model is
reliable when the program hands it the information and the narrow tool
surface the request needs, and unreliable otherwise
{[}\texttt{evidence/reports/d0}--\texttt{d5}, \texttt{d4}{]}.

{\def\LTcaptype{none} % do not increment counter
\begin{longtable}[]{@{}
  >{\raggedright\arraybackslash}p{(\linewidth - 4\tabcolsep) * \real{0.3333}}
  >{\raggedright\arraybackslash}p{(\linewidth - 4\tabcolsep) * \real{0.3333}}
  >{\raggedright\arraybackslash}p{(\linewidth - 4\tabcolsep) * \real{0.3333}}@{}}
\toprule\noalign{}
\begin{minipage}[b]{\linewidth}\raggedright
Request type
\end{minipage} & \begin{minipage}[b]{\linewidth}\raggedright
Model chooses among \textasciitilde25 tools
\end{minipage} & \begin{minipage}[b]{\linewidth}\raggedright
Program supplies information / acts directly
\end{minipage} \\
\midrule\noalign{}
\endhead
\bottomrule\noalign{}
\endlastfoot
Edit an open document & invents paths (\texttt{/Users/yourname/…}),
≈2/10 & server attaches the real path (read from the window's
accessibility attribute), doc tools only: \textbf{11/11} \\
Open an app (``微信'', ``vscode'') & 5 of 9 replies claim ``opened''
with no tool call & resolver over localized app names + server fast
path: \textbf{20/20} \\
Today's schedule & --- & small helper app with its own permission text
(Python is denied silently; AppleScript does not expand recurring
events): \textbf{5/5} test events, one brief per day \\
Unread mail / summarize message \emph{N} & --- & Mail scripting
(read-only) + local summary: \textbf{10/10} summaries without factual
errors \\
Website content & old video/search heuristics answer instead & one
read-only OpenCLI tool and the site's commands: Bilibili 3/3, Zhihu 3/3
(after login) \\
WeChat unread & --- & the app exposes six accessibility nodes; window
screenshot + on-device OCR: names 10/10, times 10/10, unread markers
4/4 \\
WeChat send & --- & search → click the exact-name entry → read the chat
title back → paste → check the new bubble: 7/7 delivered \\
\end{longtable}
}

Every step of the WeChat sender is followed by a check, because GUI
effects are not predictable: in the first live send, Return in the
search box opened a web search instead of the contact, and the title
check stopped the send before anything was typed. An unknown or misspelt
recipient is looked up first (``wenbo yen'' → ``Wenbo yan?'', sent only
after confirmation); a name that is not in the contacts is never sent
to.

\textbf{Two near-misses.} (i) Asked to edit a file that did not exist,
the server fell back to listing every open document; the model then ran
three replacements on the user's real source file in another editor
(each replaced a line with itself, so the content was unchanged). The
fix is least privilege rather than instructions: each turn may write
only the files the request names and that exist, never a substitute, and
no-op edits are rejected (re-test 11/11, the real file untouched). (ii)
A bare ``确认'' (confirm) with no message pending reached the model,
which replied ``sent to Wenbo yan'' --- nothing had been sent.
Confirmations are now answered by the program, and the per-turn usage
log marks ``claimed done without any tool call'' (including ``已发给'',
\emph{sent}) as its own status; this status is also how the false
``opened'' replies were caught.

\subsubsection{7.2 Choices to a decision
model}\label{choices-to-a-decision-model}

Which handler a request needs was first decided by hand-written
patterns. They are fast but know only the phrasings written into them
and misfire on look-alikes. We replaced them with one decision-model
question per turn (``which of these twelve kinds of request is this?'',
plus the day for schedule questions and read-or-summarize for mail),
sent at the start of the turn so it overlaps the server's other checks.
Below 0.8 confidence the request goes to the general local-model loop;
if the decision model is unavailable, the patterns decide. Handlers
still cut out text parameters themselves (recipient, message number, app
name); when a WeChat request has no explicit boundary between recipient
and message, the local model extracts them and the recipient must occur
verbatim in the request {[}\texttt{evidence/reports/d7},
\texttt{eval/intent\_v1}{]}:

{\def\LTcaptype{none} % do not increment counter
\begin{longtable}[]{@{}
  >{\raggedright\arraybackslash}p{(\linewidth - 8\tabcolsep) * \real{0.2000}}
  >{\raggedright\arraybackslash}p{(\linewidth - 8\tabcolsep) * \real{0.2000}}
  >{\raggedright\arraybackslash}p{(\linewidth - 8\tabcolsep) * \real{0.2000}}
  >{\raggedright\arraybackslash}p{(\linewidth - 8\tabcolsep) * \real{0.2000}}
  >{\raggedright\arraybackslash}p{(\linewidth - 8\tabcolsep) * \real{0.2000}}@{}}
\toprule\noalign{}
\begin{minipage}[b]{\linewidth}\raggedright
Phrasing set
\end{minipage} & \begin{minipage}[b]{\linewidth}\raggedright
Classes
\end{minipage} & \begin{minipage}[b]{\linewidth}\raggedright
Patterns
\end{minipage} & \begin{minipage}[b]{\linewidth}\raggedright
Decision model
\end{minipage} & \begin{minipage}[b]{\linewidth}\raggedright
With threshold (unsure → local model)
\end{minipage} \\
\midrule\noalign{}
\endhead
\bottomrule\noalign{}
\endlastfoot
held-out 1 (30, run once) & 10 & 15/30 & \textbf{30/30} & 30/30 \\
held-out 2 (30, run once) & 12 & 16/30 & \textbf{30/30} & 28/30 (two
correct answers below 0.8) \\
\end{longtable}
}

Pattern errors include ``明天会下雨吗'' (will it rain tomorrow) →
schedule, ``B站是哪一年成立的'' (when was Bilibili founded) → website,
and ``打开第一封看看'' (open the first one, i.e.~e-mail) → open an app;
most new phrasings were simply not recognised. The same split serves
website requests: the decision model picks the site command, row count
and search term (from candidates the program cuts out of the request),
and whether the user only wants the list. On ten held-out phrasings it
chose the command correctly 10/10; a regex shortcut tried first got 2 of
7 wrong {[}\texttt{evidence/reports/d5/clef-routing}{]}. Asking about a
concrete item (``does the request ask to summarise, compare, pick the
highest or filter?'') worked better than asking for a judgement (``is a
plain list enough?''): the list-only decision was wrong on 2 of 12
development phrasings with confidences of 0.6--0.8, and on 1 of 12 with
confidences mostly above 0.83 --- the same pattern as in our
decision-model benchmark, where whole-state judgements scored 0.64 and
per-item extraction 1.00
{[}\texttt{evidence/reports/decision-thx01/report.md}{]}.

\textbf{Latency.} On the M4 laptop the 4B model reads about 300 prompt
tokens/s and writes about 26 tokens/s, so a reply's cost is set by
prompt length and answer length. Per-stage timings in the usage log
located the slow paths:

{\def\LTcaptype{none} % do not increment counter
\begin{longtable}[]{@{}
  >{\raggedright\arraybackslash}p{(\linewidth - 6\tabcolsep) * \real{0.2500}}
  >{\raggedright\arraybackslash}p{(\linewidth - 6\tabcolsep) * \real{0.2500}}
  >{\raggedright\arraybackslash}p{(\linewidth - 6\tabcolsep) * \real{0.2500}}
  >{\raggedright\arraybackslash}p{(\linewidth - 6\tabcolsep) * \real{0.2500}}@{}}
\toprule\noalign{}
\begin{minipage}[b]{\linewidth}\raggedright
Request
\end{minipage} & \begin{minipage}[b]{\linewidth}\raggedright
Before
\end{minipage} & \begin{minipage}[b]{\linewidth}\raggedright
After
\end{minipage} & \begin{minipage}[b]{\linewidth}\raggedright
Cause
\end{minipage} \\
\midrule\noalign{}
\endhead
\bottomrule\noalign{}
\endlastfoot
明天会下雨吗 (rain tomorrow?) & 36.5 s, web search & 0.8 s & weather
path required the word ``天气''; now routed \\
B站是哪一年成立的 (knowledge) & 16 s, ``cannot find'' & 2.3 s, ``2009''
& tool loop with a 2,200--3,600-token prompt → short fixed chat
prompt \\
知乎热榜前三 (top three on Zhihu) & 14 s & 4.5 s; 0.4 s repeated & model
choosing the command (4.2 s) and retelling the list (5.7 s) → decision
model + program listing; 5-minute cache for public lists \\
知乎上大家在聊什么 (what is Zhihu discussing) & 27 s & 4.8 s & model
re-chose the command and retold ten rows (353 tokens) → run the decision
model's first choice, model answers in three sentences \\
\end{longtable}
}

The routing call costs about US\$0.00013 at list price; 248 calls in
this study stayed within the free allocation. A live demonstration
recorded on the user's Mac is in
\texttt{docs/assets/demo/kage-live-demo.mp4}.

\subsection{8. Related and concurrent
work}\label{related-and-concurrent-work}

\textbf{Learning from experience in context.} Skill libraries (Voyager;
Wang et al., 2023), experience memory and extracted insights (ExpeL;
Zhao et al., 2024; Agent Workflow Memory; Wang et al., 2025), evolving
playbooks (ACE; Zhang et al., 2025) and skills co-trained with RL
(SkillRL; Xia et al., 2026). These mostly use strong models; we find
that a 4B student retrieves such guidance reliably but often does not
act on it (§4) unless asked to ground it (§5.3).

\textbf{Self-improvement of the harness or the weights.} Harness search
(ADAS; Hu et al., 2024; Meta-Harness; Lee et al., 2026; Harness²,
concurrent, Google, 2026) changes prompts, tools and memory around a
model; weight-level methods include STaR (Zelikman et al., 2022),
context distillation (Snell et al., 2022), expert iteration (Anthony et
al., 2017) and, concurrently, Terminal-RSI (Yang et al., 2026), whose
recursive self-rewrite (RSR) re-executes tasks with runbooks before SFT
and RL at 27B. Our consolidation loop borrows RSR's guided re-execution
explicitly; the in-context learning and adoption findings predate our
reading of it (worklog timestamps). Gradient-based data selection (LESS;
Xia et al., 2024) is the replay selector we found counter-productive for
forgetting (§6); replay data came from AgentNet (OpenCUA; Wang et al.,
2025) and Toucan (Xu et al., 2025).

\textbf{Interfaces and routing.} SWE-agent (Yang et al., 2024) is the
closest work on agent-computer interfaces (§5); the ARTEMIS Android
automation project prompted our interface ablation. Routing a request to
a handler with a classifier is classic intent detection; LLM routers
such as RouteLLM (Ong et al., 2025) route between \emph{models} for
cost, and semantic routers match embeddings against example utterances.
Our router asks a calibrated choice model one question per turn and
keeps a confidence fallback to the general loop (§7.2). A longer survey
is in \texttt{docs/report/related-work.md}.

\subsection{9. Limitations}\label{limitations}

Small samples (18--90 runs per arm; several p-values are directional);
one student model; procedurally generated tasks with templated
instructions (the TF-IDF gate relies on this); desktop tasks are few and
partly read-only for safety; the pre-save checker's goals were extracted
once and cached; consolidation used a single LoRA configuration and no
RL. Desktop runs were on one Mac with one llama-server; GPU and CPU
non-determinism means single-run comparisons are unreliable, which is
why every comparison here is same-session and replicated. The daily-use
results (§7) come from one user on one Mac with tens of requests per
type; the phrasing sets were written by the authors, not collected from
users. Routing sends each request's text to a hosted model; mail and
WeChat content, documents and screenshots stay on the machine, but a
privacy-strict setting would need a local router (the cached
decision-model labels are a natural distillation set).

\subsection{10. Conclusion}\label{conclusion}

For a small local agent, ``it retrieved the right guidance and still
failed'' is common, and the cause is often not capability. In order:
make the observation and tools state what the student needs without
inference or escaping (one field, one tool), make it ground teacher
rules in the current task before acting, and only then consolidate ---
measuring held-out behaviour, because the samples that teach a family
are the ones that overwrite general skills. On a US\$4 budget this took
a 4B model to 41/45 on real macOS applications. In daily use the same
lesson goes one step further: let a decision model make the choices and
the program take the actions, and keep the small model for reading and
writing what only it can.

\subsection{References}\label{references}

\begin{itemize}
\tightlist
\item
  Anthony, T., Tian, Z., Barber, D. (2017). Thinking Fast and Slow with
  Deep Learning and Tree Search. NeurIPS. arXiv:1705.08439.
\item
  Biderman, D., Portes, J., et al.~(2024). LoRA Learns Less and Forgets
  Less. TMLR. arXiv:2405.09673.
\item
  Google (2026). Harness² (concurrent work). arXiv:2610.14224.
\item
  Google. ARTEMIS: Android automation. github.com/google/artemis.
\item
  Hu, S., Lu, C., Clune, J. (2024). Automated Design of Agentic Systems.
  arXiv:2408.08435.
\item
  Lee, Nair, Zhang, Lee, Khattab, Finn (2026). Meta-Harness.
  arXiv:2603.28052.
\item
  Ong, I., et al.~(2025). RouteLLM: Learning to Route LLMs from
  Preference Data. ICLR. arXiv:2406.18665.
\item
  Snell, C., Klein, D., Zhong, R. (2022). Learning by Distilling
  Context. arXiv:2209.15189.
\item
  Wang, G., et al.~(2023). Voyager: An Open-Ended Embodied Agent with
  Large Language Models. arXiv:2305.16291.
\item
  Wang, X., et al.~(2025). OpenCUA: Open Foundations for Computer-Use
  Agents (AgentNet). NeurIPS. arXiv:2508.09123.
\item
  Wang, Z. Z., Mao, J., Fried, D., Neubig, G. (2025). Agent Workflow
  Memory. ICML. arXiv:2409.07429.
\item
  Xia, M., Malladi, S., Gururangan, S., Arora, S., Chen, D. (2024).
  LESS: Selecting Influential Data for Targeted Instruction Tuning.
  ICML. arXiv:2402.04333.
\item
  Xia, P., Yao, H., et al.~(2026). SkillRL: Evolving Agents via
  Recursive Skill-Augmented Reinforcement Learning. arXiv:2602.08234.
\item
  Xu, Z., et al.~(2025). TOUCAN: Synthesizing 1.5M Tool-Agentic Data
  from Real-World MCP Environments. arXiv:2510.01179.
\item
  Yang, J., Jimenez, C. E., Wettig, A., Lieret, K., Yao, S., Narasimhan,
  K., Press, O. (2024). SWE-agent: Agent-Computer Interfaces Enable
  Automated Software Engineering. NeurIPS. arXiv:2405.15793.
\item
  Yang et al.~(2026). Terminal-RSI: Recursive Self-Improving Terminal
  Agents via Recursive Self-Rewrite and Policy-Refine. arXiv:2610.22192.
\item
  Zelikman, E., Wu, Y., Mu, J., Goodman, N. (2022). STaR: Bootstrapping
  Reasoning With Reasoning. NeurIPS. arXiv:2203.14465.
\item
  Zhang, Q., et al.~(2025). Agentic Context Engineering: Evolving
  Contexts for Self-Improving Language Models. arXiv:2510.04618.
\item
  Zhao, A., et al.~(2024). ExpeL: LLM Agents Are Experiential Learners.
  AAAI. arXiv:2308.10144.
\end{itemize}

\begin{center}\rule{0.5\linewidth}{0.5pt}\end{center}

\emph{Evidence index}: E-1/E-2 \texttt{evidence/reports/e1*},
\texttt{e2*}; C2 \texttt{c2/}; C3 \texttt{c3/}; R2 \texttt{r2/}; R3--R3e
\texttt{r3/}, \texttt{r3b/}--\texttt{r3e/}; C4 \texttt{c4/analysis.md},
\texttt{c4/round2.md}; G0 \texttt{g0/analysis.md}; daily use D0--D7
\texttt{d0/}--\texttt{d7/} (decision routing \texttt{d5/clef-routing/},
\texttt{d7/}); raw Kaggle runs
\texttt{evidence/kaggle/v*-results.tar.gz}.

\end{document}
